\chapter{Duże liczby kardynalne: nieosiągalność i mierzalność}
\label{ch:large-cardinals}
\chaptermark{Kardynały nieosiągalne i mierzalne}

\section{Od niezależności zdań do silniejszych aksjomatów}

W Rozdziałach~\ref{ch:constructible}, \ref{ch:boolean-valued-independence}
i~\ref{ch:symmetric-choice} badaliśmy, jak różne modele teorii
mnogości rozstrzygają CH i aksjomat wyboru. Duże liczby kardynalne
dostarczają innego rodzaju informacji: postulują istnienie obiektów,
których własności wymuszają bogatą strukturę całego uniwersum zbiorów.

Pracujemy w ZFC jako metateorii, chyba że wyraźnie zaznaczymy
inaczej. Wszystkie występujące tutaj kardynały utożsamiamy
z początkowymi liczbami porządkowymi. Symbol $V_\alpha$
oznacza $\alpha$-ty etap hierarchii kumulatywnej, wprowadzony
w Rozdziale~\ref{ch:zf-axioms}.
Nie zakładamy istnienia żadnego nieosiągalnego ani mierzalnego
kardynału: twierdzenia postaci „jeżeli istnieje $\kappa$, to ...”
są implikacjami udowadnianymi w ZFC.

\section{Regularność, silna graniczność i nieosiągalność}

\begin{definicja}[Współkońcowość kardynału]
\label{def:lc37-cofinality}
Niech $\kappa$ będzie nieskończoną liczbą porządkową.
Jej \emph{współkońcowością} $\operatorname{cf}(\kappa)$ nazywamy
najmniejszą liczbę porządkową $\delta$, dla której istnieje
ściśle rosnąca funkcja $f:\delta\to\kappa$ o obrazie
nieograniczonym w $\kappa$.
Kardynał $\kappa$ jest \emph{regularny},
gdy $\operatorname{cf}(\kappa)=\kappa$,
a \emph{singularny}, gdy $\operatorname{cf}(\kappa)<\kappa$.
\end{definicja}

\begin{definicja}[Silna graniczność i kardynały nieosiągalne]
\label{def:lc37-inaccessible}
Kardynał nieskończony $\kappa$ jest \emph{silnie graniczny},
gdy $2^\lambda<\kappa$ dla każdego kardynału $\lambda<\kappa$.
Kardynał $\kappa$ jest \emph{słabo nieosiągalny},
gdy jest nieprzeliczalny, regularny i jest granicznym
kardynałem. Jest \emph{silnie nieosiągalny},
gdy jest nieprzeliczalny, regularny i silnie graniczny.
Bez dodatkowego przymiotnika przez
\emph{kardynał nieosiągalny} rozumiemy w tym rozdziale
kardynał silnie nieosiągalny.
\end{definicja}

\begin{uwaga}[Dlaczego z definicji wyłączamy $\omega$?]
\label{rem:lc37-omega}
Kardynał $\omega$ jest regularny i spełnia
$2^n<\omega$ dla każdej skończonej liczby $n$.
Dlatego bez warunku nieprzeliczalności sam $\omega$
byłby silnie graniczny i regularny, choć nie
jest dużym kardynałem. Również warunek „graniczny”
nie implikuje silnej graniczności:
zapisujemy oba pojęcia osobno.
\end{uwaga}

\begin{lemat}[Kardynał nieosiągalny jest graniczny]
\label{lem:lc37-limit}
Jeżeli $\kappa$ jest silnie nieosiągalny,
to nie jest następnikiem żadnego kardynału.
W szczególności $\kappa$ jest graniczną liczbą porządkową.
\end{lemat}
\begin{proof}
Gdyby $\kappa=\lambda^+$, to twierdzenie Cantora
dawałoby $\lambda^+\le 2^\lambda$, podczas gdy
silna graniczność wymagałaby $2^\lambda<\kappa=\lambda^+$.
Otrzymujemy sprzeczność.
Nieprzeliczalny kardynał, który nie jest następnikiem
kardynału, jest granicznym kardynałem, a więc
graniczną liczbą porządkową.
\end{proof}

\begin{lemat}[Rozmiary etapów hierarchii poniżej $\kappa$]
\label{lem:lc37-small-ranks}
Jeżeli $\kappa$ jest silnie nieosiągalny,
to $|V_\alpha|<\kappa$ dla każdej liczby porządkowej
$\alpha<\kappa$.
W konsekwencji $|x|<\kappa$ dla każdego $x\in V_\kappa$.
\end{lemat}
\begin{proof}
Indukujemy pozaskończenie po $\alpha<\kappa$.
Dla $\alpha=0$ teza jest oczywista.
Jeżeli $|V_\alpha|<\kappa$, to
$V_{\alpha+1}=\mathcal P(V_\alpha)$ i
$|V_{\alpha+1}|=2^{|V_\alpha|}<\kappa$.

Niech $0<\lambda<\kappa$ będzie graniczną liczbą
porządkową. Mamy
$V_\lambda=\bigcup_{\alpha<\lambda}V_\alpha$.
Rodzina ma moc $|\lambda|<\kappa$, a każdy
jej wyraz ma moc mniejszą od $\kappa$.
Ponieważ $\kappa$ jest regularny, supremum
tych mocy jest pewnym $\mu<\kappa$.
Zatem
$|V_\lambda|\le|\lambda|\cdot\mu<\kappa$,
przy czym dla mocy skończonych nierówność
wynika tym bardziej bezpośrednio.
Jeśli $x\in V_\kappa$, to $x\in V_\alpha$
dla pewnego $\alpha<\kappa$ i
$|x|\le|V_\alpha|<\kappa$.
\end{proof}

\begin{twierdzenie}[Nieosiągalny etap jest modelem ZFC]
\label{thm:lc37-rank-model}
Niech $\kappa$ będzie kardynałem silnie nieosiągalnym
w uniwersum spełniającym ZFC. Wówczas zbiorowa,
tranzytywna struktura $(V_\kappa,\in)$ spełnia ZFC.
\end{twierdzenie}
\begin{proof}
Wszystkie elementy elementów $V_\kappa$ nadal
należą do $V_\kappa$, stąd tranzytywność.
Ekstensjonalność i ufundowanie są dziedziczone
z uniwersum. Pusty zbiór i $\omega$ należą do
$V_\kappa$, gdyż $\kappa>\omega$ jest graniczny.

Jeżeli $x,y\in V_\kappa$, to ich para i suma
mają rangę nie większą niż pewna skończona
liczba następczych kroków ponad
$\max\{\operatorname{rank}(x),\operatorname{rank}(y)\}$.
Z graniczności $\kappa$ wynika, że para i suma
należą do $V_\kappa$.
Tak samo $\mathcal P(x)\in V_\kappa$ dla
$x\in V_\kappa$.

Jeżeli formuła $\varphi(x,\vec p)$ z parametrami
$\vec p\in V_\kappa$ definiuje podzbiór
$a\in V_\kappa$, to relacja prawdziwości
$\varphi$ w \emph{zbiorowej strukturze} $V_\kappa$
jest definiowalna zewnętrznie.
Wyróżnianie w metateorii daje taki podzbiór $b$;
ponieważ $b\subseteq a$, jego ranga
jest nie większa od rangi $a$, a więc
$b\in V_\kappa$.

Sprawdzimy schemat zastępowania.
Niech $a,\vec p\in V_\kappa$ i załóżmy, że
$(V_\kappa,\in)$ spełnia
$\forall x\in a\ \exists!y\,\varphi(x,y,\vec p)$.
Zastępowanie w otaczającym uniwersum,
z wykorzystaniem predykatu spełniania
struktury zbiorowej, daje zbiór
$B$ wszystkich odpowiednich wartości $y\in V_\kappa$.
Z Lematu~\ref{lem:lc37-small-ranks} mamy
$|a|<\kappa$, a zatem $|B|<\kappa$.
Dla każdego $y\in B$ mamy
$\operatorname{rank}(y)<\kappa$.
Regularność daje
$\sup_{y\in B}(\operatorname{rank}(y)+1)<\kappa$;
stąd $B\in V_\kappa$.

Wreszcie dla dowolnego $a\in V_\kappa$
wybieramy w otaczającym uniwersum
relację dobrze porządkującą $a$,
korzystając z AC. Relacja ta jest
podzbiorem $a\times a$, więc jej ranga
jest ograniczona przez skończoną liczbę
następców rangi $a$. Należy zatem do
$V_\kappa$, co dowodzi AC w $V_\kappa$.
\end{proof}

\begin{wniosek}[Względna niesprzeczność z nieosiągalności]
\label{cor:lc37-inaccessible-con}
ZFC dowodzi implikacji arytmetycznej
\[
 \bigl(\exists\kappa\
    \text{$\kappa$ jest silnie nieosiągalny}\bigr)
 \ \Longrightarrow\ \operatorname{Con}(\mathrm{ZFC}).
\]
Jeżeli ZFC jest niesprzeczna, to
ZFC nie dowodzi istnienia kardynału silnie
nieosiągalnego.
\end{wniosek}
\begin{proof}
Z Twierdzenia~\ref{thm:lc37-rank-model} otrzymujemy
model zbiorowy ZFC. Formalizacja zasady poprawności
rachunku dowodowego daje wtedy
$\operatorname{Con}(\mathrm{ZFC})$;
argument jest przeprowadzalny w samej ZFC.
Jeśli ZFC dowodziłaby istnienia takiego kardynału,
dowodziłaby własnej niesprzeczności.
Przeczy to drugiemu twierdzeniu Gödla
omówionemu w Rozdziale~\ref{ch:second-incompleteness},
przy założeniu niesprzeczności ZFC.
\end{proof}

\section{Ultrafiltry zupełne i mierzalność}

\begin{definicja}[Ultrafiltr i zupełność]
\label{def:lc37-complete-ultrafilter}
Niech $I$ będzie niepustym zbiorem.
\emph{Filtrem właściwym} na $I$ nazywamy
$\mathcal U\subseteq\mathcal P(I)$ spełniający:
$I\in\mathcal U$, $\varnothing\notin\mathcal U$;
z $A,B\in\mathcal U$ wynika
$A\cap B\in\mathcal U$; z $A\in\mathcal U$
i $A\subseteq B\subseteq I$ wynika $B\in\mathcal U$.
Filtr jest \emph{ultrafiltrem}, gdy
dla każdego $A\subseteq I$ dokładnie jeden
ze zbiorów $A$, $I\setminus A$ należy do
$\mathcal U$. Jest \emph{niegłówny}, gdy
nie zawiera żadnego singletonu.

Dla nieskończonego kardynału $\kappa$ mówimy,
że filtr jest \emph{$\kappa$-zupełny}, gdy
dla każdego $\mu<\kappa$ i rodziny
$(A_\xi)_{\xi<\mu}$ jego elementów
przecięcie $\bigcap_{\xi<\mu}A_\xi$ również
należy do filtru.
\end{definicja}

\begin{definicja}[Kardynał mierzalny]
\label{def:lc37-measurable}
Nieprzeliczalny kardynał $\kappa$ jest
\emph{mierzalny}, jeśli na zbiorze $\kappa$
istnieje niegłówny $\kappa$-zupełny
ultrafiltr $\mathcal U$.
Taki ultrafiltr będziemy również nazywać
\emph{miarą zero-jedynkową} na $\kappa$;
nie chodzi tutaj o miarę liczbową o dowolnych
wartościach z przedziału $[0,1]$.
\end{definicja}

\begin{lemat}[Małe zbiory nie należą do miary]
\label{lem:lc37-uniform}
Niech $\kappa>\omega$ będzie kardynałem,
a $\mathcal U$ niegłównym $\kappa$-zupełnym
ultrafiltrem na $\kappa$. Wówczas
żaden zbiór $A\subseteq\kappa$ o mocy
$|A|<\kappa$ nie należy do $\mathcal U$.
W szczególności każdy podzbiór $\kappa$,
którego dopełnienie ma moc mniejszą od
$\kappa$, należy do $\mathcal U$.
\end{lemat}
\begin{proof}
Dla każdego $\alpha\in A$ mamy
$\kappa\setminus\{\alpha\}\in\mathcal U$.
Zupełność pozwala przeciąć $|A|<\kappa$
tych zbiorów, więc
$\kappa\setminus A\in\mathcal U$.
Właściwy filtr nie może zawierać
jednocześnie $A$ i jego dopełnienia.
Drugi wniosek jest już udowodniony.
\end{proof}

\begin{twierdzenie}[Kardynał mierzalny jest nieosiągalny]
\label{thm:lc37-measurable-inaccessible}
Jeżeli nieprzeliczalny kardynał $\kappa$
jest mierzalny, to jest regularny i
silnie graniczny, a więc silnie nieosiągalny.
\end{twierdzenie}
\begin{proof}
Niech $\mathcal U$ będzie niegłównym
$\kappa$-zupełnym ultrafiltrem na $\kappa$.

\emph{Regularność.} Gdyby
$\operatorname{cf}(\kappa)=\mu<\kappa$,
to pewien rosnący, współkońcowy ciąg
$(\gamma_\xi)_{\xi<\mu}$ dzieliłby
$\kappa$ na $\mu$ przedziałów
o mocy mniejszej od $\kappa$.
Żaden z tych przedziałów nie należałby
do $\mathcal U$ na mocy
Lematu~\ref{lem:lc37-uniform}.
Ich dopełnienia należałyby więc do
$\mathcal U$, a przecięcie tych
$\mu<\kappa$ dopełnień byłoby puste.
To przeczy $\kappa$-zupełności
i właściwości filtru.

\emph{Silna graniczność.} Załóżmy,
że istnieje kardynał $\lambda<\kappa$
taki, że $2^\lambda\ge\kappa$.
W ZFC istnieje wtedy iniekcja
$f:\kappa\to\mathcal P(\lambda)$.
Dla każdego $\xi<\lambda$ ultrafiltr
wybiera dokładnie jeden z dwóch zbiorów
\[
 A_\xi=\{\alpha<\kappa:\xi\in f(\alpha)\},
 \qquad \kappa\setminus A_\xi.
\]
Niech $B_\xi$ oznacza wybrany element
$\mathcal U$. Ponieważ $\lambda<\kappa$,
zbiór $B=\bigcap_{\xi<\lambda}B_\xi$
należy do $\mathcal U$.
Każde dwie liczby $\alpha,\beta\in B$
dają jednak te same odpowiedzi
na pytania o przynależność
$\xi\in f(\alpha)$, dla wszystkich
$\xi<\lambda$. Zatem
$f(\alpha)=f(\beta)$, a iniektywność
$f$ pokazuje, że $|B|\le1$.
Jest to sprzeczne z
Lematem~\ref{lem:lc37-uniform}.
\end{proof}

\begin{przyklad}[Dlaczego nie wystarcza zwykły ultrafiltr?]
\label{ex:lc37-frechet}
Niegłówny ultrafiltr na $\omega$
można otrzymać, stosując lemat o
ultrafiltrze do filtru zbiorów
o skończonych dopełnieniach.
Jest on $\omega$-zupełny, ponieważ
przecięcia \emph{skończenie} wielu
elementów zawsze należą do filtru.
Nie jest jednak $\omega_1$-zupełny:
dla $n<\omega$ zbiory
$A_n=\{m\in\omega:m\ge n\}$
należą do takiego ultrafiltru, lecz
$\bigcap_{n<\omega}A_n=\varnothing$.
Wymóg nieprzeliczalności w definicji
mierzalności nie jest więc kosmetyczny.
\end{przyklad}

\section{Ultraprodukt uniwersum i twierdzenie Łosia}

\begin{definicja}[Klasa ultraproduktu uniwersum]
\label{def:lc37-ultrapower}
Niech $\kappa$ będzie mierzalny, a
$\mathcal U$ niegłównym $\kappa$-zupełnym
ultrafiltrem na $\kappa$. Rozważamy
klasę wszystkich funkcji zbiorowych
$f:\kappa\to V$. Określamy
\[
 f\equiv_{\mathcal U}g
 \iff\{\alpha<\kappa:f(\alpha)=g(\alpha)\}
           \in\mathcal U.
\]
Przez $[f]_{\mathcal U}$ oznaczamy
klasę równoważności (jako obiekt
klasowy), a relację przynależności
ustalamy przez
\[
 [f]_{\mathcal U}\in_{\mathcal U}[g]_{\mathcal U}
 \iff\{\alpha<\kappa:f(\alpha)\in g(\alpha)\}
           \in\mathcal U.
\]
Tak powstałą strukturę klasową
oznaczamy $\operatorname{Ult}(V,\mathcal U)$.
\end{definicja}

\begin{uwaga}[Klasy i ich reprezentanci]
\label{rem:lc37-class-ultrapower}
W ZFC nie traktujemy zbioru wszystkich
funkcji $\kappa\to V$ jako zbioru.
Zamiast ilorazu będącego zbiorem można
wybrać \emph{reprezentantów Scotta}:
spośród równoważnych funkcji wybieramy
te o najmniejszej randze i bierzemy
zbiór wszystkich takich reprezentantów.
Pozwala to kodować każdy element
ultraproduktu jednym zbiorem
i rozważać definiowalną klasę
tych kodów.
\end{uwaga}

\begin{twierdzenie}[Twierdzenie Łosia dla ultraproduktu $V$]
\label{thm:lc37-los}
Niech $\kappa$ będzie mierzalny i
$\mathcal U$ jak w
Definicji~\ref{def:lc37-ultrapower}.
Dla każdej ustalonej formuły
$\varphi(x_1,\ldots,x_n)$ języka
teorii mnogości oraz funkcji
$f_1,\ldots,f_n:\kappa\to V$ zachodzi
\[
 \operatorname{Ult}(V,\mathcal U)
 \models\varphi([f_1]_{\mathcal U},
                \ldots,[f_n]_{\mathcal U})
 \quad\Longleftrightarrow\quad
 \{\alpha<\kappa:
 \varphi(f_1(\alpha),\ldots,f_n(\alpha))\}
 \in\mathcal U.
\]
Prawa strona jest predykatem ZFC dla
każdej ustalonej formuły.
\end{twierdzenie}
\begin{proof}
Dla atomów $x=y$ i $x\in y$
równoważność jest definicją obu
relacji. Przy koniunkcji używamy
zamknięcia $\mathcal U$ na przecięcia,
a przy negacji własności ultrafiltru:
dla dowolnego podzbioru $\kappa$
należy doń dokładnie jeden spośród
tego zbioru i jego dopełnienia.

Rozważmy kwantyfikator egzystencjalny.
Jeżeli klasa $[g]_{\mathcal U}$
jest świadkiem formuły, to hipoteza
indukcyjna daje zbiór dużych indeksów,
na których $g(\alpha)$ jest świadkiem
dla argumentów $f_i(\alpha)$.
W drugą stronę niech
\[
 D=\{\alpha<\kappa:
  \exists y\ \psi(y,f_1(\alpha),\ldots,
                         f_n(\alpha))\}\in\mathcal U.
\]
Dla każdego $\alpha\in D$ wybieramy
jeden zbiór $g(\alpha)$ spełniający
formułę; dla $\alpha\notin D$
przyjmujemy $g(\alpha)=\varnothing$.
Wybór jest poprawny w ZFC: najpierw
wybieramy najmniejszą możliwą rangę
świadka, następnie jednego ze świadków
w odpowiednim zbiorze tej rangi,
stosując AC do rodziny indeksowanej
\emph{zbiorem} $D$.
Otrzymujemy funkcję zbiorową
$g:\kappa\to V$. Zbiór indeksów,
dla których $\psi(g(\alpha),\vec f(\alpha))$
jest prawdziwa, zawiera $D$,
a więc należy do $\mathcal U$.
Hipoteza indukcyjna daje świadka $[g]$
po lewej stronie. To zamyka indukcję
po budowie ustalonej formuły.
\end{proof}

\begin{lemat}[Dobre ufundowanie ultraproduktu]
\label{lem:lc37-wellfounded}
Przy założeniach
Twierdzenia~\ref{thm:lc37-los}
relacja $\in_{\mathcal U}$ jest
dobrze ufundowana w sensie
zewnętrznym: nie istnieje ciąg
$([f_m]_{\mathcal U})_{m<\omega}$
spełniający
$[f_{m+1}]_{\mathcal U}\in_{\mathcal U}
 [f_m]_{\mathcal U}$ dla wszystkich $m$.
\end{lemat}
\begin{proof}
Gdyby taki ciąg istniał,
dla każdego $m$ mielibyśmy
\[
 A_m=\{\alpha<\kappa:
      f_{m+1}(\alpha)\in f_m(\alpha)\}
      \in\mathcal U.
\]
Ponieważ $\omega<\kappa$, z zupełności
$\mathcal U$ wynikałoby
$\bigcap_{m<\omega}A_m\in\mathcal U$.
Wybierając $\alpha$ z tego
niepustego przecięcia,
dostalibyśmy nieskończony ciąg
$f_0(\alpha)\ni f_1(\alpha)\ni
f_2(\alpha)\ni\cdots$, sprzeczny
z aksjomatem ufundowania w V.
\end{proof}

\begin{twierdzenie}[Zanurzenie elementarne z miary]
\label{thm:lc37-embedding}
Jeżeli $\kappa$ jest mierzalny,
to istnieje definiowalny z parametru
$\mathcal U$ tranzytywny model klasowy
$M\models\mathrm{ZFC}$ oraz
zanurzenie elementarne $j:V\to M$
takie, że
\[
 j(\alpha)=\alpha\quad(\alpha<\kappa),
 \qquad j(\kappa)>\kappa.
\]
W szczególności najmniejsza liczba
porządkowa przesuwana przez $j$,
zwana \emph{punktem krytycznym},
jest równa $\kappa$.
\end{twierdzenie}
\begin{proof}
Niech $\mathcal U$ będzie miarą na
$\kappa$. Dzięki Twierdzeniu~\ref{thm:lc37-los}
ultraprodukt jest elementarnym obrazem
uniwersum przy odwzorowaniu
$x\mapsto[c_x]_{\mathcal U}$, gdzie
$c_x(\alpha)=x$.
Z Lematu~\ref{lem:lc37-wellfounded}
relacja członkostwa jest dobrze ufundowana
dla ciągów długości $\omega$.
Jest również ekstensjonalna
z elementarności, a ponadto
\emph{zbiorolokalna}: dla ustalonej
$f$ członkowie $[f]$ mają reprezentantów
$g$ takich, że
$g(\alpha)\in f(\alpha)$ na odpowiednim
zbiorze indeksów. Można więc ograniczyć
zakres wszystkich takich $g$ do zbioru
$\bigl(\bigcup_{\alpha<\kappa}f(\alpha)\bigr)
\cup\{\varnothing\}$,
zmieniając wartości poza zbiorem
miary jeden. Zbiór odpowiednich funkcji
jest zbiorem. Istnieje zatem
kolaps Mostowskiego tej dobrze
ufundowanej, ekstensjonalnej,
zbiorolokalnej klasy definiowalnej
do tranzytywnej klasy $M$.
Złożenie odwzorowania stałych funkcji
z kolapsem oznaczamy przez $j$.
Twierdzenie Łosia daje elementarność,
a więc $M\models\mathrm{ZFC}$.

Pokażemy indukcyjnie, że
$j(\alpha)=\alpha$ dla $\alpha<\kappa$.
Dla $\alpha=0$ teza jest oczywista.
Jeśli $j(\beta)=\beta$ dla $\beta<\alpha$,
to dowolny element $j(\alpha)$
jest reprezentowany funkcją
$f:\kappa\to\alpha$ na zbiorze
miary jeden. Podzielenie tego zbioru
na mniej niż $\kappa$ włókien
$f^{-1}(\{\beta\})$, $\beta<\alpha$,
i $\kappa$-zupełność ultrafiltru
dają jedno włókno należące do
$\mathcal U$. Zatem $[f]=[c_\beta]$
dla pewnej $\beta<\alpha$,
co po kolapsie oznacza, że każdy
element $j(\alpha)$ jest mniejszy
od $\alpha$. Ponieważ
$j(\beta)\in j(\alpha)$ dla
wszystkich $\beta<\alpha$,
otrzymujemy $j(\alpha)=\alpha$.

Niech $i:\kappa\to\kappa$
będzie funkcją identycznościową.
Dla każdego $\alpha<\kappa$
zbiór $\{\xi:\alpha<\xi\}$ ma
dopełnienie mocy mniejszej od
$\kappa$, więc należy do $\mathcal U$
na mocy Lematu~\ref{lem:lc37-uniform}.
Stąd w uporządkowaniu ultraproduktu
\[
 [c_\alpha]_{\mathcal U}
 <[i]_{\mathcal U}
 <[c_\kappa]_{\mathcal U}.
\]
Po kolapsie dostajemy
$\alpha<[i]$ dla wszystkich
$\alpha<\kappa$ oraz
$[i]<j(\kappa)$. Zatem
$\kappa\le[i]<j(\kappa)$,
czyli $j(\kappa)>\kappa$.
\end{proof}

\section{Miary normalne i diagonalne przecięcia}

\begin{definicja}[Normalny ultrafiltr]
\label{def:lc37-normal}
Niech $\kappa$ będzie nieprzeliczalnym
kardynałem regularnym, a $\mathcal U$
niegłównym $\kappa$-zupełnym
ultrafiltrem na $\kappa$.
Mówimy, że $\mathcal U$ jest
\emph{normalny}, gdy dla każdego
$A\in\mathcal U$ i funkcji
$f:A\to\kappa$ spełniającej
$f(\alpha)<\alpha$ dla
wszystkich $0<\alpha\in A$,
istnieją $B\subseteq A$,
$B\in\mathcal U$ oraz $\xi<\kappa$
takie, że $f|_B$ jest stale
równa $\xi$.
Warunek w punkcie $0$ nie ma
znaczenia, ponieważ singleton
$\{0\}$ nie należy do $\mathcal U$.
\end{definicja}

\begin{definicja}[Przecięcie diagonalne]
\label{def:lc37-diagonal}
Dla ciągu $(A_\xi)_{\xi<\kappa}$
podzbiorów $\kappa$ definiujemy
\emph{przecięcie diagonalne}
\[
 \mathop{\triangle}_{\xi<\kappa}A_\xi
 =\{\alpha<\kappa:
    \forall\xi<\alpha\ (\alpha\in A_\xi)\}.
\]
\end{definicja}

\begin{twierdzenie}[Charakterystyka normalności]
\label{thm:lc37-normal-diagonal}
Niech $\kappa$ będzie nieprzeliczalnym
regularnym kardynałem, a $\mathcal U$
niegłównym $\kappa$-zupełnym
ultrafiltrem na $\kappa$.
Wówczas następujące warunki są równoważne:
\begin{enumerate}
\item $\mathcal U$ jest normalny;
\item dla każdej rodziny
$(A_\xi)_{\xi<\kappa}\subseteq\mathcal U$
jej przecięcie diagonalne należy
do $\mathcal U$.
\end{enumerate}
\end{twierdzenie}
\begin{proof}
Przyjmijmy normalność i niech
$D=\mathop{\triangle}_{\xi<\kappa}A_\xi$.
Jeżeli $D\notin\mathcal U$, to
$C=\kappa\setminus D\in\mathcal U$.
Dla $\alpha\in C\setminus\{0\}$
niech $f(\alpha)$ będzie
najmniejszym $\xi<\alpha$
takim, że $\alpha\notin A_\xi$.
Funkcja jest regresywna.
Z normalności dostajemy
$B\in\mathcal U$, $B\subseteq C$,
na którym $f$ jest stale równa
pewnej $\xi<\kappa$.
Wtedy $B\cap A_\xi=\varnothing$,
co jest sprzeczne z tym, że
oba zbiory należą do ultrafiltru.

Odwrotnie, załóżmy własność
przecięć diagonalnych.
Niech $A\in\mathcal U$,
a $f:A\setminus\{0\}\to\kappa$
będzie regresywna.
Przypuśćmy, że żadne włókno
$f^{-1}(\{\xi\})$ nie należy
do $\mathcal U$.
Dla każdej $\xi<\kappa$
zbiór
$A_\xi=\kappa\setminus f^{-1}(\{\xi\})$
należy do $\mathcal U$.
Wtedy $A\cap
 \mathop{\triangle}_{\xi<\kappa}A_\xi$
należy do $\mathcal U$,
a więc jest niepusty.
Dla $\alpha$ z tego przecięcia,
$\alpha>0$, mamy $f(\alpha)<\alpha$
oraz $\alpha\notin
 f^{-1}(\{f(\alpha)\})$, sprzeczność.
Jedno z włókien należy więc
do $\mathcal U$, co daje normalność.
\end{proof}

\begin{twierdzenie}[Kardynał mierzalny ma miarę normalną]
\label{thm:lc37-normal-exists}
Dla każdego kardynału mierzalnego
$\kappa$ istnieje niegłówny
$\kappa$-zupełny normalny ultrafiltr
na $\kappa$.
\end{twierdzenie}
\begin{proof}
Niech $\mathcal U$ będzie miarą,
a $j:V\to M$ zanurzeniem z
Twierdzenia~\ref{thm:lc37-embedding}.
Definiujemy
\[
 \mathcal W=\{X\subseteq\kappa:
                  \kappa\in j(X)\}.
\]
Jest to zbiór, a nie klasa
właściwa: funkcja $j$ jest
definiowalna z ustalonego parametru
$\mathcal U$, więc działa
wyodrębnianie na $\mathcal P(\kappa)$.

Z elementarności $j$ mamy
$j(\kappa\setminus X)
 =j(\kappa)\setminus j(X)$
i $j(X\cap Y)=j(X)\cap j(Y)$.
Ponieważ $\kappa<j(\kappa)$,
dokładnie jeden spośród $X$,
$\kappa\setminus X$ należy do
$\mathcal W$, a przecięcia i
nadzbiory zachowują przynależność.
Dla $\alpha<\kappa$ mamy
$j(\{\alpha\})=\{\alpha\}$,
zatem $\{\alpha\}\notin\mathcal W$.

Jeżeli $\mu<\kappa$ i
$(X_\xi)_{\xi<\mu}$ jest rodziną
elementów $\mathcal W$, to
$j(\mu)=\mu$ i elementarność
daje
\[
 j\left(\bigcap_{\xi<\mu}X_\xi\right)
  =\bigcap_{\xi<\mu}j(X_\xi).
\]
Wszystkie prawe zbiory zawierają
$\kappa$, więc przecięcie
należy do $\mathcal W$.
To dowodzi $\kappa$-zupełności.

Na koniec niech $A\in\mathcal W$,
a $f:A\to\kappa$ będzie regresywna
poza punktem $0$.
Ponieważ $\kappa\in j(A)$ i
$\kappa>0$, elementarność daje
$j(f)(\kappa)<\kappa$.
Jest to więc pewna
$\xi<\kappa$, a $j(\xi)=\xi$.
Zatem
\[
 \kappa\in j(\{\alpha\in A:
                          f(\alpha)=\xi\}),
\]
co znaczy, że włókno
$f^{-1}(\{\xi\})$ należy do
$\mathcal W$. Ultrafiltr
$\mathcal W$ jest normalny.
\end{proof}

\section{Hierarchia założeń i granice otrzymanych wyników}

\begin{twierdzenie}[Wynikające implikacje]
\label{thm:lc37-hierarchy}
W ZFC prawdziwe są następujące twierdzenia:
\begin{enumerate}
\item jeśli istnieje mierzalny kardynał,
to istnieje silnie nieosiągalny kardynał;
\item jeśli istnieje silnie nieosiągalny
kardynał, to istnieje tranzytywny
model zbiorowy ZFC;
\item jeśli istnieje mierzalny kardynał,
to istnieje nietrywialne zanurzenie
elementarne $j:V\to M$, dla pewnego
tranzytywnego, definiowalnego
modelu klasowego $M$ ZFC.
\end{enumerate}
Żadnej z odwrotnych implikacji nie
stwierdzamy w tym rozdziale.
\end{twierdzenie}
\begin{proof}
Pierwsza część jest
Twierdzeniem~\ref{thm:lc37-measurable-inaccessible},
druga wynika z
Twierdzenia~\ref{thm:lc37-rank-model},
a trzecia z
Twierdzenia~\ref{thm:lc37-embedding}.
\end{proof}

\begin{uwaga}[Co jest nowym aksjomatem?]
\label{rem:lc37-new-axioms}
ZFC nie zapewnia istnienia
silnie nieosiągalnego ani mierzalnego
kardynału. W odróżnieniu od CH i AC,
których niezależność omawialiśmy
wcześniej, tutaj rozważamy
\emph{wzmocnienia} teorii aksjomatami
istnienia dużych kardynałów.
Drugie twierdzenie Gödla wyklucza
dowód istnienia nieosiągalnego w
niesprzecznej ZFC, lecz samo
to rozumowanie nie daje dowodu
względnej niesprzeczności
takiego dodatkowego aksjomatu.
\end{uwaga}

\begin{uwaga}[Dalsze zagadnienia]
\label{rem:lc37-next}
Słaba zwartość kardynałów,
liczby kardynalne Mahla,
silna zwartość, superzwartość
oraz rozbudowane twierdzenia o
zanurzeniach elementarnych
wymagają odrębnych definicji
i dowodów. W tym rozdziale
nie zakładamy ich własności.
Nie utożsamiamy także
ultrafiltrów mierzalnych z
miarami Lebesgue'a.
\end{uwaga}

\begin{uwaga}[Zakres kontroli formalnej]
\label{rem:lc37-formal}
Sprawdzenie skończonych przykładów
ultrafiltrów i kompilacja LaTeX
nie są formalizacją twierdzeń
o dużych kardynałach w Lean.
Taką weryfikację należy
raportować osobno na podstawie
rzeczywistego przebiegu kompilatora.
\end{uwaga}
